\section{Boundaries}

This slice verifies a single ballot-polling outcome assertion under sampling
assumptions. BRAVO confirms that a declared stop-or-continue transcript follows
from reported two-candidate totals, a risk limit, and an ordered list of ballot
observations. It does not verify ballot custody, voter eligibility, or any
ballot-level CVR comparison claim; those live in other RCOUNT contracts.

The statistical boundary is also deliberate. The implemented recurrence is the
classic two-candidate BRAVO likelihood ratio in integer parts-per-million. It is
not Minerva, ALPHA, Kaplan-Markov, or a stratified or comparison risk measure,
and it does not decide whether a jurisdiction's legal RLA method permits
replacement sampling, round-by-round planning, or multi-candidate assertions.
Multi-candidate contests must be reduced to the relevant pairwise assertions
before this contract applies.

Finally, BRAVO replay assumes the sample order is already accounted for. It
checks that the observed marks support the declared decision, not that the
ballots were the ones a public seed should have selected; that reproducibility
claim is the separate sampler-replay contract of V.7. As with every RCOUNT audit
slice, the public-evidence privacy boundary remains with the jurisdiction:
publishing per-ballot observations can expose rare ballot styles, so small-batch
publication rules are an external responsibility.
